% Generated draft: requires mathematical and source review.
\documentclass[11pt]{article}
\usepackage[margin=25mm]{geometry}
\usepackage{fontspec}
\setmainfont{Latin Modern Roman}
\usepackage{amsmath,amssymb,mathtools,unicode-math}
\setmathfont{Latin Modern Math}
\usepackage{enumitem,needspace,xurl,hyperref}
\setlength{\parindent}{0pt}
\setlength{\parskip}{5pt}
\setlength{\emergencystretch}{3em}
\newcommand{\N}{\mathbb{N}}
\newcommand{\Z}{\mathbb{Z}}
\newcommand{\Q}{\mathbb{Q}}
\newcommand{\R}{\mathbb{R}}
\newcommand{\C}{\mathbb{C}}
\newcommand{\F}{\mathbb{F}}
\newcommand{\A}{\mathbb{A}}
\newcommand{\PP}{\mathbb{P}}
\newcommand{\E}{\mathbb{E}}
\newcommand{\eps}{\varepsilon}
\newcommand{\dd}{\,\mathrm{d}}
\begin{document}
\section*{UnsolvedMath ID: 800937}
\textbf{Source problem code:} catalogue entry
\subsection{Definitions and mathematical statement}
Let \(M\) be a compact smooth manifold with smooth boundary \(N=\partial M\), and let \(g\) be a Riemannian metric on \(M\). For \(p,q\in N\), define the boundary distance by
\[
d_g(p,q)=\inf_{\substack{\gamma:[0,1]\to M\ \text{piecewise }C^1\\ \gamma(0)=p,\ \gamma(1)=q}}\int_0^1\sqrt{g_{\gamma(t)}(\dot\gamma(t),\dot\gamma(t))}\,\dd t.
\]
The source describes this function as finite and nonnegative. If no such curve exists, the infimum is \(+\infty\); thus the stated finite-valued convention requires an additional connectivity qualification. No connectedness or nonemptiness assumption on \(N\) is explicitly supplied.

The source's definition contains a self-comparison: it says that a competing metric \(g'\) must be isometric to \(g'\). To obtain a nontrivial rigidity question, the apparent intended comparison is between \((M,g)\) and \((M,g')\). With this correction, an unmarked interpretation is the following property:
\[
\begin{gathered}
\text{For every Riemannian metric }g'\text{ on }M,\\
\left[d_{g'}(p,q)=d_g(p,q)\text{ for all }p,q\in N\right]
\Longrightarrow
\left[\text{there exists a diffeomorphism }\Phi:M\to M
\text{ with }\Phi^*g'=g\right].
\end{gathered}
\]
A stronger, boundary-marked interpretation requires the same conclusion with the additional condition \(\Phi|_N=\operatorname{id}_N\). The supplied wording does not specify whether the isometry must fix the boundary pointwise, so these interpretations must remain distinguished.

\textbf{Question 17 (Boundary rigidity problem).} What conditions on \(M\), \(N\), and \(g\) imply boundary rigidity in the intended comparison sense? The target is a request for sufficient conditions, not an assertion that every such metric is rigid. The self-comparison correction, the boundary-marking convention, and the finite-distance qualification require human review.
\subsection{Short English statement}
Which conditions ensure that distances between boundary points, measured through the manifold, determine its Riemannian metric up to isometry? The supplied definition apparently contains a self-comparison typo and does not specify whether the isometry must fix the boundary pointwise. Its finite-distance convention also needs a connectivity qualification.
\subsection{Sources}
\begin{itemize}
\item Local UnsolvedMath registry, permanent ID \(800937\), exact source problem code \texttt{catalogue entry}, in the AMR Open Problem Lists collection. The supplied record labels the problem \emph{open}; this is a registry label, not an independently verified literature assessment.
\item Burns and Matveev, \emph{Open Problems and Questions about Geodesics} (supplied year: 2019), Question 17, with supplied attribution/context: Boundary rigidity problem. Supplied source URL: \url{https://arxiv.org/abs/1308.5417}.
\item Supplied provenance reports access on \texttt{2026-10-16} and extraction by \texttt{arxiv\_latex}. These are supplied provenance details; no independent source visit or verification is claimed.
\end{itemize}
\end{document}
